% PUBLIC / SANITIZED RESUME -- served at joshuacowell.com/resume/
% Do NOT add employer names, program names, or classified duty detail here.
% Full-detail versions live in resume/source/ and are gitignored.
\documentclass[10pt,letterpaper]{article}

% ============================================
% PACKAGES
% ============================================
\usepackage[margin=0.4in, top=0.4in, bottom=0.4in]{geometry}
\usepackage[T1]{fontenc}
% Times-like serif font widely used on professional/defense resumes
\usepackage{mathptmx}
\usepackage{enumitem}
\usepackage{titlesec}
\usepackage{xcolor}
\usepackage{hyperref}
\usepackage{tabularx}
\usepackage{setspace}

% ============================================
% COLORS - Professional Navy for Defense
% ============================================
\definecolor{navyblue}{RGB}{0, 40, 85}
\definecolor{darkgray}{RGB}{50, 50, 50}
\definecolor{medgray}{RGB}{100, 100, 100}

% ============================================
% HYPERLINK SETUP
% ============================================
\hypersetup{
    colorlinks=false,
    pdfborder={0 0 0}
}

% ============================================
% REMOVE PAGE NUMBERS
% ============================================
\pagenumbering{gobble}

% ============================================
% SECTION FORMATTING
% ============================================
\titleformat{\section}{\normalsize\bfseries}{}{0em}{\textcolor{navyblue}}[\textcolor{navyblue}{\titlerule}]
\titlespacing{\section}{0pt}{8pt}{4pt}

% Uppercase section titles manually
\renewcommand{\section}[1]{
    \vspace{5pt}
    {\normalsize\bfseries\textcolor{navyblue}{\MakeUppercase{#1}}}
    \vspace{0pt}
    
    {\color{navyblue}\hrule height 0.5pt}
    \vspace{3pt}
}

% ============================================
% SPACING
% ============================================
\setlength{\parindent}{0pt}
\setlength{\parskip}{0pt}
\setlist[itemize]{leftmargin=0.15in, topsep=1pt, itemsep=0pt, parsep=0pt, label=\textcolor{navyblue}{\textbullet}}

% ============================================
% CUSTOM COMMANDS
% ============================================
\newcommand{\resumeEntry}[4]{
    \textbf{#1} \hfill \textcolor{medgray}{#2} \\
    \textit{\textcolor{darkgray}{#3}} \ifx&#4&\else\hfill \textcolor{medgray}{#4}\fi
    \vspace{0pt}
}

\newcommand{\projectEntry}[1]{
    \textbf{#1}
    \vspace{0pt}
}

% ============================================
% DOCUMENT
% ============================================
\begin{document}

% ============================================
% HEADER
% ============================================
\begin{center}
    {\fontsize{20}{24}\selectfont\bfseries\textcolor{navyblue}{JOSHUA COWELL}}

    \vspace{3pt}

    {\small\textbf{\textcolor{navyblue}{Active Top Secret / SCI Clearance}} \enspace\textcolor{medgray}{$|$}\enspace \textbf{\textcolor{navyblue}{Full-Scope Polygraph}}}

    \vspace{3pt}

    {\small\color{darkgray}
    Lynchburg, VA \enspace\textcolor{medgray}{$|$}\enspace
    \href{mailto:joshuacowell2005@gmail.com}{joshuacowell2005@gmail.com} \enspace\textcolor{medgray}{$|$}\enspace
    \href{https://joshuacowell.com}{joshuacowell.com}
    }
\end{center}

\vspace{0pt}

% ============================================
% EDUCATION
% ============================================
\section{Education}

\textbf{Bachelor of Science in Mechanical Engineering} \hfill \textcolor{medgray}{Expected May 2027} \\
\textit{\textcolor{darkgray}{Liberty University, Lynchburg, VA}} \hfill \textcolor{medgray}{GPA: 3.85}

\vspace{1pt}
{\small \textit{Selected Coursework:} Dynamics, Fluid Dynamics, Heat Transfer, Materials Engineering, Differential Equations, Thermal-Fluids Laboratory}

\vspace{1pt}
{\small \textit{Founding Member:} American Nuclear Society (ANS) student chapter (2025); AIAA student chapter (2026)}

\vspace{1pt}
{\small \textit{Leadership:} Liberty Rocketry -- Chairman, Board of Advisors (Jul 2026 -- Present)}

\vspace{1pt}

% ============================================
% ENGINEERING EXPERIENCE
% ============================================
\section{Engineering Experience}

\resumeEntry{Advanced Weapons Systems Analyst Intern}{May 2026 -- Present}{U.S. Government}{}
\begin{itemize}
    \item Conduct physical testing of engineering systems and analyze results to calibrate simulation models
    \item Developed a Python-based geospatial analysis tool to meet team requirements, using AI-assisted development tooling to iterate on the design
    \item Perform technical assessments of advanced engineering systems alongside multidisciplinary teams, translating findings into decision-ready summaries
    \item Brief senior leadership and author concise written products that communicate complex technical data clearly
    \item \textit{Additional detail available upon request}
\end{itemize}

\vspace{1pt}

\resumeEntry{Chief Systems Engineer and Team Lead}{Jun 2025 -- Jul 2026}{Liberty Rocketry}{}

\vspace{1pt}
{\small \textit{\textcolor{darkgray}{Earlier roles: Assistant Chief Engineer; Senior Structural Engineer; Fin Fluid-Dynamics Engineer}}}
\begin{itemize}
    \item Placed 18th of 141 international collegiate teams at the 2026 IREC
    \item Led a 60-member student team in the design, testing, and launch readiness of a high-powered rocket
    \item Oversaw engineering decisions across propulsion, avionics, recovery, and aerodynamics subsystems
    \item Facilitated communication among internal teams, external vendors, and university faculty to improve project execution
    \item Reviewed SolidWorks, ANSYS, MATLAB, OpenRocket, and RocketPy models and simulations used for team design decisions
\end{itemize}

\vspace{1pt}

\resumeEntry{Thermal Hydraulics Engineering Intern}{May 2025 -- Dec 2025}{Framatome Nuclear}{}
\begin{itemize}
    \item Developed and validated a Python predictive-wear model against documented component wear, reducing prediction error by 75\% versus the legacy model
    \item Built a Python/xarray algorithm processing six-dimensional datasets to calculate safe operating limits, reducing analysis time from 800+ hours to minutes
    \item Performed system-level thermal-hydraulic and FIV analyses using AFT Fathom and custom tools to support component qualification
    \item Authored technical documentation and presented findings to engineering leadership; received a \$2,500 technical-presentation award
\end{itemize}

\vspace{1pt}

\resumeEntry{Lead Undergraduate Research Assistant}{Sep 2024 -- Present}{Liberty University}{}
\begin{itemize}
    \item Lead a five-person team conducting split-Hopkinson pressure bar (SHPB) testing for dynamic material characterization
    \item Develop test procedures and fixtures; use PicoScope strain data to generate stress-strain curves and correlate high-speed imaging with SEM results
\end{itemize}

\vspace{1pt}

% ============================================
% TECHNICAL SKILLS
% ============================================
\section{Technical Skills}

\begin{tabularx}{\textwidth}{@{}>{}l X@{}}
\textbf{Systems \& Test:} & System integration, test planning, model calibration, subsystem coordination, technical documentation \\[2pt]
\textbf{Analysis \& CAE:} & ANSYS Fluent, ANSYS Mechanical, AFT Fathom, OpenRocket, RocketPy, CFD, FEA \\[2pt]
\textbf{CAD \& Manufacturing:} & SolidWorks, Onshape, GD\&T, technical drawings, 3D printing, CNC machining \\[2pt]
\textbf{Programming \& Lab:} & Python (NumPy, pandas, xarray), MATLAB, SHPB testing, SEM, high-speed imaging \\
\end{tabularx}

\vspace{1pt}

% ============================================
% SELECTED PROJECTS
% ============================================
\section{Selected Projects}

\projectEntry{Unmanned Aerial Vehicle Design \& Analysis}

{\small \textit{\textcolor{darkgray}{Tools: SolidWorks, ANSYS Fluent, 3D Printing}}}
\begin{itemize}
    \item Designed and integrated a modular UAV airframe, propulsion system, and payload using replaceable 3D-printed components; evaluated performance with CFD
\end{itemize}

\vspace{1pt}

\projectEntry{3D-Printed Model Rocket}

{\small \textit{\textcolor{darkgray}{Tools: ANSYS Fluent, OpenRocket, Additive Manufacturing}}}
\begin{itemize}
    \item Engineered a reusable threaded airframe with a hybrid recovery system; simulations predicted Mach 0.5+ and approximately 3,000 ft altitude
\end{itemize}

\end{document}
